\section{Exact dynamics of training-pattern moments}
\label{sec:pattern-quotient}
The residual convention matters. In this section define
$\rho_n=f_\theta(x_n)-x_n$ and
$A_g=N^{-1}\sum_n g_n\rho_nx_n^T$.
If instead the residual is $x_n-f_\theta(x_n)$, all three displayed minus
signs below become plus signs. Both conventions describe the same flow.

\begin{lemma}[Fixed-pattern moment dynamics]\label{lem:pattern-quotient}
On an open time interval where every training preactivation is nonzero
and the neuronwise training patterns are fixed, let
\[
 I_g=\{i:\mathbf1_{u_i^Tx_n>0}=g_n\text{ for all }n\},\quad
 M_g=\sum_{i\in I_g}w_iu_i^T,\quad
 U_g=\sum_{i\in I_g}u_iu_i^T,\quad
 W_g=\sum_{i\in I_g}w_iw_i^T.
\]
Then the exact closed equations are
\begin{align*}
 \dot M_g&=-A_gU_g-W_gA_g,\\
 \dot U_g&=-A_g^TM_g-M_g^TA_g,\\
 \dot W_g&=-A_gM_g^T-M_gA_g^T,
 \qquad f(x_n)=\sum_g g_nM_gx_n.
\end{align*}
All moments include their entire groups; no concentration or low-rank
approximation is assumed.
\end{lemma}
\begin{proof}
Direct differentiation of the stated normalized loss gives
$\dot w_i=-A_gu_i$ and $\dot u_i=-A_g^Tw_i$ for $i\in I_g$.
Differentiate each outer product and sum. For example
$\dot w_iu_i^T+w_i\dot u_i^T=-A_gu_iu_i^T-w_iw_i^TA_g$.
The other two identities follow in the same way, with the transposes shown.
The network identity follows from $(u_i^Tx_n)_+=g_nu_i^Tx_n$.
Consequently $A_g$ depends only on the fixed patterns, the data, and the
moments, proving closure on this interval.
\end{proof}

Put
\[
 K_g=\begin{pmatrix}U_g&M_g^T\\M_g&W_g\end{pmatrix},\qquad
 B_g=\begin{pmatrix}0&A_g^T\\A_g&0\end{pmatrix}.
\]
Then $K_g\succeq0$, $\dot K_g=-B_gK_g-K_gB_g$ and
$\operatorname{tr}(U_g-W_g)$ is constant on the fixed-pattern interval.
The same equations hold for a partition refined by fixed cohort membership
and fixed probe signs, since every subgroup with training pattern $g$
uses the same $A_g$.

\paragraph{What the quotient retains and loses.}
The training-pattern moments determine the training outputs on that cell.
They do not, by themselves, determine a new probe's output, even if every
individual probe sign is specified. For a concrete example use a single
training vector $e_1$, $a,b>0$, and two balanced neurons $w_i=u_i$.
Compare
\[
 u_1=(a,b),\quad u_2=(a,-b)
 \quad\hbox{with}\quad
 u'_1=(7a/5,b/5),\quad u'_2=(a/5,-7b/5).
\]
This is an orthogonal mixing of the two neuron columns. Both training
patterns are active, and in both states the $-e_2$ probe is inactive on
neuron 1 and active on neuron 2. All three training-group moments are
identical, equal to $\operatorname{diag}(2a^2,2b^2)$. Yet
\[
 f(-e_2)=(ab,-b^2),\qquad
 f'(-e_2)=(7ab/25,-49b^2/25).
\]
Thus probe-refined moments (or an equivalent signed probe observable)
are necessary for an OOD interface. This example concerns the sufficiency
of proposed observable data; it is not a counterexample to canonical
initialized reversal.

\paragraph{Gate-event transfers.}
Let $y_i=(u_i,w_i)$ and $K_i=y_iy_i^T$. At an isolated switch of neuron $i$
from training group $g$ to $g'$, the underlying state is continuous, but
its group representation changes by
\[
 K_g^+=K_g^- -K_i,\qquad K_{g'}^+=K_{g'}^-+K_i.
\]
Simultaneous switches sum their individual contributions. The group moment
does not identify $K_i$. Indeed different rank-one decompositions, including
the orthogonal mixing above, have the same $K_g$ and different $K_i$.
An estimate of the form
\begin{equation}\label{eq:invalid-transfer-bound}
 \|K_i-\widehat K_i\|\le C\|K_g-\widehat K_g\|
\end{equation}
therefore fails for every finite universal $C$. A transfer enclosure must
retain the boundary neuron's rank-one moment or other sufficient geometry.
This is an information loss of the unaugmented quotient, not a failure of
the proposed augmented quotient/event strategy.

\paragraph{Reference census.}
The deterministic census of the stored numerical reference has 53 distinct
training patterns at $t=0$, 121 at $t=1$, 140 at $t=2$, 134 at $t=3.4$,
and 134 at $t=4$. A 29-neuron inactive training group persists at these
reference knots. Over 20000 adjacent reference-knot pairs there are 185211
sample--neuron sign differences. These are diagnostic counts, not certified
actual-flow events; multiple crossings between knots could be missed.
At entry a 16-neuron training-pattern group mixes strong-cohort membership.
It must be split by membership to recover the repository $Q_G$ from group
moments. The stepwise CSV also records all group sizes and a nonrigorous
linear-motion near-boundary count. No census number is a proof premise.
